\documentclass[12pt]{article}
\usepackage[utf8]{inputenc}
\usepackage[T1]{fontenc}
\usepackage{hyperref}
\title{Kawaiisu Language Resources for K-12 Curriculum Development}
\author{Compiled by Claude Sonnet 5.5 (claude-sonnet-5-5) for Andr\'as Kornai}
\date{}
\begin{document}
\maketitle

\section{Introduction}
Kawaiisu (ISO 639-3: xaw; Glottocode: kawa1283) is a Numic language of the Uto-Aztecan family, from the Tehachapi area of California. It is critically endangered, with very few elderly speakers. No school programme is known.

\section{Primary Educational Resources}
The Kawaiisu Language \& Cultural Center offers language classes and DVDs for home learning (as of 2012). No downloadable textbook or primer was found: \url{https://actaonline.org/kawaiisu-language-cultural-center}

\section{Linguistic and Reference Materials}
\begin{itemize}
\item Zigmond, Booth and Munro, \emph{Kawaiisu: A Grammar and Dictionary, With Texts}, UC Publications in Linguistics 119 (University of California Press, 1991): \url{https://mitpressbookstore.mit.edu/book/9780520097476}
\item The same book at another bookseller: \url{https://page158books.com/book/9780520097476}
\item NSF documentation project BCS-1160669, which added about 500 dictionary entries along with recordings and transcriptions: \url{https://grantome.com/grant/NSF/BCS-1160669}
\item Kawaiisu language (Wikipedia): \url{https://en.wikipedia.org/wiki/Kawaiisu_language}
\end{itemize}

\section{Religious and Community Resources}
The Kawaiisu Language \& Cultural Center (see above) is the community organisation found in the searches. Its own website is listed as possibly unreliable, so check it directly. Local history: \url{https://www.tehachapinews.com/opinion/community-voices-the-origins-of-tehachapi/article_f29fce07-aacf-4b9c-b166-66221a9328b9.html}

\section{Digital and Audio Resources}
The NSF project above produced recordings and transcriptions. I did not find where they are hosted. The searches found no standalone online dictionary.

\section{Curriculum Development Context}
The Zigmond--Booth--Munro grammar and dictionary is the main written base. The Center's classes and DVDs are the only known teaching materials. Any curriculum would be a revitalization effort built on archival documentation and the few remaining speakers.
\end{document}